\begin{recipe}
[% 
    preparationtime = {\unit[30]{min}},
    portion = {\portion{2}},
    source = {\protect\recipesource{https://www.orodiparma.de/italienische-rezepte/spaghetti-carbonara}{ORO di Parma: Spaghetti Carbonara}},
    calory = {660kcal | 26g Protein | 83g KH | 25g Fett},
    price = {2,20€},
    created = {11.06.2026},
    %rating = {1},
]
{Vegetarische Carbonara mit getrockneten Tomaten}

%\graph{% pictures
%    big=pic/hauptgericht/19_vegetarische-carbonara-getrocknete-tomaten
%}

\introduction{%
Carbonara-Style ohne Fleisch, dafür mit getrockneten Tomaten als salzig-intensiver Umami-Part. Durch Ei, Eigelb, Parmesan und Pastawasser wird die Sauce cremig, ohne dass man Sahne reinschummeln muss.
}

\ingredients{
    200 g & Spaghetti\index{Getreide!Pasta}\\
    80 g & Getrocknete Tomaten in Öl\index{Gemüse!Tomate}\\
    1 & Ei\index{Eier!Ei}\\
    1 & Eigelb\index{Eier!Ei}\\
    40 g & Parmesan\index{Milchprodukte \& Alternativen!Parmesan}\\
    1 Zehe & Knoblauch\\
    1 EL & Olivenöl oder Öl der getrockneten Tomaten\\
    & Schwarzer Pfeffer\\
    & Salz\\
    Optional & Petersilie
}

\preparation{%
    \step%
    Parmesan fein reiben. Ei, Eigelb, Parmesan und ordentlich frisch gemahlenen schwarzen Pfeffer in einer Schüssel verrühren.
    
    \step%
    Getrocknete Tomaten abtropfen lassen und in Streifen schneiden. Knoblauch fein hacken.
    
    \step%
    Spaghetti in gut gesalzenem Wasser al dente kochen. Vor dem Abgießen unbedingt etwas Pastawasser auffangen.
    
    \step%
    Währenddessen Olivenöl oder etwas Öl aus dem Tomatenglas in einer Pfanne erhitzen. Getrocknete Tomaten darin 2--3 Minuten anbraten, bis sie leicht dunkler werden.
    
    \step%
    Knoblauch zugeben und kurz mitbraten, bis er duftet. Nicht verbrennen lassen, sonst schmeckt es direkt nach Fehler.
    
    \step%
    Pfanne vom Herd ziehen. Spaghetti tropfnass dazugeben und alles gut vermengen.
    
    \step%
    Ei-Parmesan-Mischung zu den Nudeln geben und sofort gründlich durchschwenken. Nach und nach etwas Pastawasser zugeben, bis eine cremige Sauce entsteht.
    
    \step%
    Mit schwarzem Pfeffer, wenig Salz und optional Petersilie servieren.
}

\hint{
    Die Pfanne muss vom Herd, bevor die Ei-Parmesan-Mischung dazukommt. Sonst gibt es Rührei mit Nudeln, und das ist dann leider keine Carbonara mehr.
}

\suggestion[Getrocknete Tomaten]{%
    Wenn die Tomaten sehr salzig sind, beim Nudelwasser und beim finalen Abschmecken vorsichtiger sein.
}

\end{recipe}